\documentclass[UTF8]{ctexbook}
\usepackage[a4paper,margin=1in]{geometry}
\usepackage{amsmath,amssymb}
\usepackage{booktabs,longtable}
\usepackage{hyperref}
\usepackage{enumitem}
\usepackage{xcolor}
\usepackage{listings}

\hypersetup{
  colorlinks=true,
  linkcolor=blue,
  urlcolor=blue
}

\lstset{
  basicstyle=\ttfamily\small,
  frame=single,
  breaklines=true,
  columns=fullflexible
}

\setlist[itemize]{leftmargin=2em}
\setlist[enumerate]{leftmargin=2.2em}

\title{Stanford CS329A 课程补编\\把《Advanced LLM Agents》推进到《Self-Improving AI Agents》}
\author{Supplement workspace for CS294/194-280 Spring 2025 textbook}
\date{2026-04-16}

\begin{document}
\maketitle
\tableofcontents

\chapter{Stanford CS329A 课程补编：把“Advanced LLM Agents”推进到“Self-Improving AI Agents”}

\section{补编定位与证据边界}

本补编服务于现有的 Berkeley《CS294/194-280: Advanced Large Language Model Agents, Spring 2025》中文教材。它不是一门新书的完整重建，而是一个课程级扩展章：用 Stanford《CS329A: Self-Improving AI Agents, Autumn 2025》的官方课程页与官方 schedule，回答一个更强的问题：如果把 Spring 2025 Berkeley 课程看成 ``advanced LLM agents'' 的系统教材，那么 Stanford CS329A 把这条线进一步推进到了哪里？

这里必须先说清楚证据边界：
\begin{itemize}
\item 本补编的 canonical source 只有两个：\url{https://cs329a.stanford.edu/} 与 \url{https://cs329a.stanford.edu/#schedule}。
\item 我没有在 Stanford 官方课程站点上找到官方公开视频页。
\item 我没有在 Stanford 官方课程站点上找到官方 slide 页面或 lecture slide 索引。
\item 因此，本补编不是逐讲视频重建，也不是 slide 逐页讲义，而是基于官方课程概览、官方 schedule 结构、官方 reading 标题和课程组织方式做出的教材化扩展。
\item 存在一个公开 YouTube playlist，但它更适合被描述为 instructor-affiliated，而不是 Stanford 课程官网明确公布的官方 source-of-truth，所以这里只做说明，不把它当作覆盖依据。
\end{itemize}

这意味着本补编有一个非常重要的方法论约束：凡是超出官方课程页和 reading 标题字面信息的解释，都应被理解为教材化推断或结构化解读，而不是对未公开 lecture 内容的伪装复原。

\section{为什么 CS329A 是现有 Berkeley 教材的强补充}

Berkeley Spring 2025 的主线，是把 advanced LLM agents 拆成几个核心子系统：推理时计算、学习推理、记忆与规划、coding agents、多模态 agents、formal mathematics、abstraction/discovery，以及安全与权限边界。

Stanford CS329A 的官方课程概览则把这些对象重新组织成一个更强的系统目标：self-improvement。课程页没有只说 agent 很强，而是明确列出一条自我改进链路：constitutional AI、verifiers、scaling test-time compute、combining search with LLMs、train-time scaling with RL、tool use、code、memory、multimodal interaction、multi-step reasoning and planning、evaluation frameworks、coding agents、research assistants in STEM，以及 robotics。

如果把这条链路教材化，可以写成下面这个抽象循环：
\[
\pi_{t+1} = \mathcal{U}(\pi_t, \mathcal{T}, \mathcal{V}, \mathcal{M}, \mathcal{S}, \mathcal{D})
\]

这里的写法不是 Stanford 官网原文，而是对官方课程概览的教材化抽象：
\begin{itemize}
\item $\pi_t$：第 $t$ 轮 agent 系统。
\item $\mathcal{T}$：tools 与 code execution 反馈。
\item $\mathcal{V}$：verifiers 与外部检查器。
\item $\mathcal{M}$：memory 与长期上下文结构。
\item $\mathcal{S}$：search、planning 与 test-time compute。
\item $\mathcal{D}$：训练数据、交互轨迹与评测反馈。
\item $\mathcal{U}$：把这些反馈汇总成系统更新的机制。
\end{itemize}

Berkeley Spring 2025 更像是在解释这些部件各自如何工作；Stanford Autumn 2025 则更像是在问：怎样把这些部件串成一个持续改进的 agent loop。这正是它对现有教材的第一个关键补强。

\section{第一条扩展主线：从“推理技巧”升级到“自我改进回路”}

\subsection{Test-time compute 不再只是推理技巧，而是自我改进的前半段}

Stanford schedule 把 2025 年 9 月 26 日直接设成 ``Test-time Compute Scaling''，并配上 Large Language Monkeys、Archon、Scaling LLM Test-Time Compute Optimally can be More Effective than Scaling Model Parameters，以及 How Do Large Language Monkeys Get Their Power (Laws)? 这些 reading。

只看标题就足够说明课程意图：这不是把 repeated sampling 或 search 当作零散 trick，而是把它们看成一个可被系统地扩展、比较、甚至形成 scaling law 的能力来源。Berkeley 的第 1 讲已经讲了 inference-time reasoning techniques，但 Stanford 在课程结构上又往前走了一步：它把 test-time compute 直接放进 self-improving agent 的总框架中。这意味着推理时计算不只是为了在一次任务上答对，而是为了给后续的 selection、verification、feedback learning 和 system redesign 提供更高质量的候选轨迹。

\subsection{Robust Verification 把“验证器”提前变成中心对象}

2025 年 9 月 29 日的主题是 ``Robust Verification''。相应 reading 包括 Shrinking the Generation-Verification Gap with Weak Verifiers、Training Verifiers to Solve Math Word Problems、Let's Verify Step by Step，以及 Math-Shepherd。

这和 Berkeley Spring 2025 的一个差异非常关键。Berkeley 的 verifier 主要出现在推理、formal mathematics 和 theorem proving 这些部分；而 Stanford 在课表编排上更像是在说：没有 verifier，就谈不上 reliable self-improvement。

这会直接改变教材中的逻辑顺序。原书里，verifier 往往被读者理解为某些高难数学或形式化任务的附加组件；但从 CS329A 的编排看，verifier 更像是所有 agent 改进循环里的通用裁判。只要 agent 要从自己的候选行为中学习，generation-verification gap 就会成为首要问题。

\subsection{Tool/code feedback 和 planning 被合并进同一个闭环}

2025 年 10 月 3 日是 ``Learning from feedback with tools/code''，10 月 6 日是 ``Multi-step Reasoning/Planning''。这两次课连在一起看，比单独看更重要。

官方 reading 标题给出的信号是：ReAct 说明 reasoning 与 acting 要耦合；RLEF 说明 code execution feedback 可以成为强化学习信号；Constitutional AI 说明 feedback 不一定来自人，也可以来自规则化的 AI feedback；而 SWiRL、Language Agent Tree Search、SPRINT、ADaPT 与 Wider or Deeper? 则把多步推理、树搜索、分解规划、并行执行都拉进来。

因此，在 Stanford 的课程设计里，tool use 不是 ``会不会调 API'' 这种孤立技能，planning 也不是 ``让模型多想一步'' 的装饰。二者被放在一起，是因为有工具的 environment 才会产生足够强的反馈，而多步规划又决定 agent 如何消费这些反馈。

把它写成一个教材化伪代码，更接近 CS329A 的课程精神：
\begin{lstlisting}
Initialize agent policy and tool interface
Generate multi-step plans under a limited compute budget
Execute tool/code actions and collect environment feedback
Use verifiers or rule-based critics to filter candidate trajectories
Store useful traces, failures, and reusable subplans
Update the agent or search policy for the next round
\end{lstlisting}

这段伪代码不是官网原文，但它忠实反映了官方 overview 与 schedule 的组合含义：planning、tool use、feedback learning、memory 与 verification 被组织成一个统一的 self-improvement workflow。

\subsection{Train-time scaling/Scaling RL 补上自我改进的后半段}

2025 年 10 月 10 日的主题是 ``Train Time Scaling/Scaling RL''，reading 包括 STaR、DeepSeekMath 和 DAPO。

Berkeley Spring 2025 里，读者已经见过 DPO/PPO、reasoning post-training 和 formal RL 在特定 lecture 中的作用。Stanford 这里的增量不是 ``又多几篇 RL 论文''，而是把 train-time scaling 明确写成自我改进回路的另一半：前半段是 test-time compute、search、planning、verification，后半段则是如何把这些交互轨迹回流为训练信号。

\section{第二条扩展主线：从 discovery 章节走向 open-ended research automation}

\subsection{Open-Ended Evolution 把发现问题本身推到前台}

2025 年 10 月 13 日的主题是 ``Open-Ended Evolution of Self-Improving Agents''。官方 reading 标题包括 Automated design of agentic systems、The AI Scientist: Towards Fully Automated Open-Ended Scientific Discovery，以及 AlphaEvolve: A Gemini-powered coding agent for designing advanced algorithms。

Berkeley Spring 2025 的第 11 讲已经讨论了 abstraction 和 discovery，但 Stanford 在课程命名上把 ``open-ended evolution'' 直接放进主干。这个差别很大：Berkeley 更像在问 agent 如何帮助 theorem proving、abstraction 和 scientific discovery；Stanford 更像在问 agent 能否持续地设计新的 agent system、新的 research workflow，甚至新的 algorithmic artifact。

\subsection{Search \& Deep Research Agents 把 coding 和 research synthesis 连到一起}

2025 年 10 月 17 日的主题是 ``Self improvement with Search \& Deep Research Agents''，reading 包括 Competition-Level Code Generation with AlphaCode、AlphaCode 2 Technical Report，以及 Search-o1: Agentic Search-Enhanced Large Reasoning Models。

单看标题，就可以看出 Stanford 想把两件原本常被分开的事情放在一起：代码与程序搜索，以及 research-style 搜索与综合。这对现有 Berkeley 教材的意义是：原书已经有 coding agents、formal mathematics、abstraction/discovery 三条线，但 Stanford 的补充会迫使我们把它们读成一条连续谱：代码生成是可执行搜索，theorem proving 是 verifier-backed 搜索，deep research 则是资料、假设与证据之间的开放式搜索。

\section{第三条扩展主线：把 agent infrastructure 讲得更像系统工程}

\subsection{从 chatbots 到 agents：post-training 被重新定义}

2025 年 10 月 20 日的 guest lecture 题目是 ``Evolution of Post-training from Chatbots to Agents''。虽然官方 schedule 没有给 reading，但这个标题本身就很值得写进补编。它说明 Stanford 把 post-training 看成 agent emergence 的制度性条件，而不只是对话模型礼貌化或 alignment 的尾声。

因为官方站点没有公开 lecture 细节，这里不能虚构具体方法；但课程题目已经足够说明其教学方向：agentic behavior 不是只靠 inference-time prompting 长出来的，post-training 本身也在发生从 chatbot 到 agent 的范式迁移。

\subsection{Agentic Frameworks for Software Engineering 让 coding agent 更接近“生产系统”}

2025 年 11 月 3 日的主题是 ``Agentic Frameworks for Software Engineering''，reading 包括 CodeMonkeys、KernelBench，以及 Improving Parallel Program Performance with LLM Optimizers via Agent-System Interfaces。

这组标题直接把 Berkeley 第 5 讲的 coding agent 讨论往前推了一层。Spring 2025 里，我们已经看到 coding agents 与 vulnerability detection；但 Stanford 这里强调的是：software engineering 是 test-time compute 的落地场景，benchmark 不该只看 pass@k，还要看性能、系统接口和并行程序优化，agent 不是单一模型，而是 model + interface + runtime + evaluator 的整体。

\subsection{Memory 被提升为长期系统能力，而不是局部增强技巧}

2025 年 11 月 7 日的 guest lecture 是 ``Augmenting Agents with Memory''，reading 包括 Cartridges、MemGPT，以及 CacheBlend。

这与 Berkeley 第 3 讲形成直接呼应。Spring 2025 的教材里，memory 已经出现，但 Stanford 的课表说明 memory 现在更像一个完整基础设施层：它关系到长期任务分解、上下文复用与缓存，也关系到 agent 是否能把过去的交互沉淀成未来的能力。

\section{第四条扩展主线：把 evaluation 放到 agent loop 的中轴}

\subsection{late-term reasoning guest lectures 的信号}

2025 年 11 月 10 日的 guest lecture 是 ``LLM Reasoning''，11 月 14 日是 ``Towards AI Superhuman Reasoning: AlphaProof, AlphaGeometry \& Gemini IMO Gold Medal''。

由于官方站点没有给出 reading，这里不能展开成独立技术综述；但它们在课程后半段的出现，已经能说明两件事：第一，Stanford 没有把 reasoning 看成开头讲完就结束的话题，而是把它反复拉回课程主干；第二，reasoning 的标尺被继续抬高到了 formal mathematics、geometry 和竞赛级表现附近。

\subsection{Agentic Evaluations \& Long-Horizon Tasks 改写了“评测”的位置}

2025 年 11 月 17 日的主题是 ``Agentic Evaluations \& Long-Horizon Tasks''，reading 包括 Measuring AI Ability to Complete Long Tasks、GDPVal，以及 DeepScholar-Bench。

这一组标题的意义非常大。Berkeley Spring 2025 当然也谈 evaluation 与 safety，但 Stanford 这里的官方安排表明：评测不再只是最后验收模型分数，评测定义了 agent 究竟在什么环境里被算作 ``改进了''，而长时程任务、经济价值任务、研究综合任务，都成为新的 primary target。

\subsection{项目制结构本身也是 evaluation philosophy 的一部分}

Stanford 官方页还清楚列出了 Homework 1、2、3，Project Proposal，Midterm presentation + Report，Final project，以及 Poster。这说明课程并不把评测局限于论文阅读或考试，而是把 ``提出问题、搭系统、做中期汇报、跑最终项目'' 本身当作学习闭环的一部分。

\section{第五条扩展主线：从多模态 agent 走向 autonomy 与 robotics}

2025 年 11 月 21 日的 guest lecture 是 ``Building Agentic Systems for Autonomy: Lessons \& Open questions''，12 月 1 日是 ``Multimodal AI Agents in Robotics''。

官方站点同样没有给出 lecture 细节，所以这里必须保持克制。但课程结构已经非常明确：Stanford 把 autonomy 与 robotics 放在课程收尾，说明它们被看成对前面所有能力的综合检验。真正进入 autonomy 或 robotics，agent 需要同时处理 perception、action、memory、planning、environment feedback、robustness，以及 long-horizon correction。这正好把 Berkeley 第 6、7 讲的多模态与 perception-to-action 主题，推进到 ``实际 autonomy systems'' 这一层。

\section{如何用这门课来重构现有 Berkeley 教材}

\begin{longtable}{p{0.23\textwidth}p{0.30\textwidth}p{0.39\textwidth}}
\toprule
现有 Berkeley 教材部分 & Stanford CS329A 的增量 & 应如何改写原书理解 \\
\midrule
\endhead
推理时计算、学习推理 & Test-time Compute Scaling、Robust Verification、Train Time Scaling/Scaling RL & 把 reasoning 写成一个从 test-time 到 train-time 的 self-improvement loop。 \\
Reasoning / Memory / Planning & Multi-step Reasoning/Planning、Augmenting Agents with Memory & 把 planning 和 memory 视为长期自我改进的运行时基础设施。 \\
Coding Agents & Learning from feedback with tools/code、Agentic Frameworks for Software Engineering、AlphaCode 2 & 从 ``会写代码'' 升级到 ``能在系统接口与执行反馈中持续优化''。 \\
Abstraction / Discovery & Open-Ended Evolution of Self-Improving Agents、AI Scientist、AlphaEvolve & 把 discovery 写成 open-ended research automation，而不仅是单点 scientific insight。 \\
Formal math / theorem proving & Robust Verification、AlphaProof 等 late-term reasoning 信号 & 把 verifier 和 formal reasoning 提前，作为通用 agent 改进机制。 \\
Multimodal agents & Multimodal AI Agents in Robotics、Autonomy & 把多模态从 perception demo 推向 embodied long-horizon autonomy。 \\
安全与评测 & Agentic Evaluations \& Long-Horizon Tasks、项目制课程结构 & 把 evaluation 变成 agent loop 的中轴，而不是章节尾声。 \\
\bottomrule
\end{longtable}

\section{本补编的使用方式}

如果你已经读完 Berkeley Spring 2025 这本书，建议用下面的顺序吸收这个 Stanford 补编：
\begin{enumerate}
\item 先重读 Berkeley 的第 1--4 讲，把 reasoning、planning、post-training 重新放进 self-improvement 框架。
\item 再回看 Berkeley 的 coding agents、formal mathematics 和 abstraction/discovery 几章，把它们读成 research automation 与 open-ended system design 的连续谱。
\item 最后把多模态 agent、安全与评测章节重新阅读，重点关注 long-horizon tasks、economic value、research synthesis 和 robotics 这几个新的评价维度。
\end{enumerate}

\section{总结}

如果说 Berkeley Spring 2025 教材回答的是 ``高级 LLM agent 由哪些关键部件构成''，那么 Stanford CS329A 进一步追问的是：当这些部件被放进长期反馈回路后，agent 如何持续改进自己？

这门课最重要的补充，不是单个新名词，而是一个课程级重心变化：reasoning 不再只是求解，verifier 不再只是附加检查器，tools 与 code 不再只是插件，memory 不再只是辅助缓存，evaluation 不再只是结果统计。它们一起变成了 self-improvement system 的组成部分。

从教材工程的角度看，这意味着现有 Berkeley 讲义项目后续如果继续扩展，最值得优先补写的，不是 ``再多几个案例''，而是把以下五条主线更清楚地前置出来：
\begin{enumerate}
\item self-improvement loop
\item verifier-centric learning
\item coding and research automation
\item long-horizon evaluation
\item autonomy and robotics as integrated end-to-end tests
\end{enumerate}

这正是 Stanford CS329A 对整本书最有价值的增量。

\end{document}
